% Generated draft: requires mathematical and source review.
\documentclass[11pt]{article}
\usepackage[margin=25mm]{geometry}
\usepackage{fontspec}
\setmainfont{Latin Modern Roman}
\usepackage{amsmath,amssymb,mathtools,unicode-math}
\setmathfont{Latin Modern Math}
\usepackage{enumitem,needspace,xurl,hyperref}
\setlength{\parindent}{0pt}
\setlength{\parskip}{5pt}
\setlength{\emergencystretch}{3em}
\newcommand{\N}{\mathbb{N}}
\newcommand{\Z}{\mathbb{Z}}
\newcommand{\Q}{\mathbb{Q}}
\newcommand{\R}{\mathbb{R}}
\newcommand{\C}{\mathbb{C}}
\newcommand{\F}{\mathbb{F}}
\newcommand{\A}{\mathbb{A}}
\newcommand{\PP}{\mathbb{P}}
\newcommand{\E}{\mathbb{E}}
\newcommand{\eps}{\varepsilon}
\newcommand{\dd}{\,\mathrm{d}}
\begin{document}
\section*{UnsolvedMath ID: 30006173}
\textbf{Source problem code:} catalogue entry
\subsection{Definitions and mathematical statement}
Under the conventional interpretation of symbolic subshifts as two-sided shifts over finite alphabets, formulate the question as follows. These conventions are not explicitly specified in the supplied record and require source confirmation.

For each integer \(k\geq 1\), let \(A_k=\{0,\ldots,k-1\}\), give \(A_k^{\Z}\) the product topology, and define the shift by
\[
(\sigma_k x)(n)=x(n+1)\qquad(n\in\Z).
\]
A subshift is a nonempty closed subset \(X\subseteq A_k^{\Z}\) satisfying \(\sigma_k[X]=X\). It is minimal if
\[
\overline{\{\sigma_k^m x:m\in\Z\}}=X
\qquad\text{for every }x\in X.
\]
Let \(\mathcal M_k\) be the space of minimal subshifts over \(A_k\), with the Borel structure inherited from the Vietoris topology on the hyperspace of nonempty compact subsets of \(A_k^{\Z}\). Use the tagged disjoint union
\[
\mathcal M=\bigsqcup_{k\geq1}\mathcal M_k
\]
as a coding space for minimal symbolic subshifts over finite alphabets.

For \(X\in\mathcal M_k\) and \(Y\in\mathcal M_\ell\), define topological conjugacy by
\[
X\mathrel{E_{\mathrm{conj}}}Y
\quad\Longleftrightarrow\quad
\text{there exists a homeomorphism }h:X\longrightarrow Y
\text{ with }h\circ(\sigma_k|_X)=(\sigma_\ell|_Y)\circ h.
\]
An equivalence relation \(E\) on a standard Borel space \(S\) is hyperfinite if there exist Borel equivalence relations \(F_j\subseteq S\times S\), indexed by \(j\in\N\), such that every \(F_j\)-class is finite and
\[
F_j\subseteq F_{j+1}\quad\text{for every }j,
\qquad E=\bigcup_{j\in\N}F_j.
\]
No uniform bound on the sizes of the finite classes is required.

A point \(x\in A_k^{\Z}\) is a Toeplitz point if
\[
\forall n\in\Z\ \exists p\in\Z_{>0}\ \forall m\in\Z,
\qquad x(n+mp)=x(n).
\]
Let \(\mathcal T\subseteq\mathcal M\) consist of those minimal subshifts that are orbit closures of Toeplitz points:
\[
X=\overline{\{\sigma_k^m x:m\in\Z\}}
\quad\text{for some Toeplitz point }x\in A_k^{\Z}.
\]
This definition includes periodic Toeplitz points; the supplied record does not state whether periodic systems are excluded by its terminology.

\textbf{Question.} Is \(E_{\mathrm{conj}}\) on \(\mathcal M\) hyperfinite? In particular, is the restricted relation
\[
E_{\mathrm{conj}}\!\upharpoonright\mathcal T
=E_{\mathrm{conj}}\cap(\mathcal T\times\mathcal T)
\]
hyperfinite, using the inherited Borel structure?

The supplied original passage mentions the minimal-subshift question and a connection with a conjecture of Thomas concerning isomorphism of complete groups. No precise formulation of that conjecture or of the connection is supplied, so neither is added as a mathematical assertion here. The Toeplitz subquestion occurs in the supplied reconstructed statement, rather than in the supplied original passage. The record labels the problem open; this is a supplied status, not an independently verified conclusion.
\subsection{Short English statement}
Does topological conjugacy of minimal symbolic subshifts admit an increasing exhaustion by Borel equivalence relations with finite classes? The supplied statement also asks this for Toeplitz subshifts. The formulation uses two-sided shifts over finite alphabets, subject to confirmation of the source conventions. This is a statement-only draft requiring human mathematical review.
\subsection{Sources}
\begin{itemize}
\item Local problem registry: \texttt{data/unsolvedmath/problems.json}, with \texttt{catalog/problems.json} as the additional registry specified by the project instructions; permanent UnsolvedMath ID \(30006173\), title \emph{Hyperfiniteness of Minimal-Subshift Conjugacy}, exact source problem code \texttt{catalogue entry}. The supplied record belongs to \emph{Oberwolfach Reports Open Problems} and gives proposed year 2025; no individual proposer is supplied.
\item Supplied primary citation: \emph{Set Theory} (2025), Oberwolfach Reports, pp.\ 75--140, DOI \texttt{10.4171/owr/2025/2}. \url{https://doi.org/10.4171/owr/2025/2}. The record describes Oberwolfach Reports as documenting mathematical workshops and their open-problem discussions.
\item Supplied literature source, \emph{Oberwolfach Report 2/2025}: \url{https://content.ems.press/assets/public/full-texts/serials/owr/22/1/14299082/online/10.4171-owr-2025-2.pdf}.
\item Supplied literature source, Deka--Peng, \emph{Generalized Oxtoby subshifts and hyperfiniteness}: \url{https://arxiv.org/abs/2502.14978}.
\item Supplied literature assessment dated \textbf{2026-08-22}: the primary report is described as stating that hyperfiniteness of conjugacy for minimal symbolic subshifts remains unknown, while 2025 work proves hyperfiniteness only for a proper generalized Oxtoby subclass. This is supplied dated curation derived from the OWR material, not an independently verified assessment. Its time-sensitive status, subclass qualification, and cited claims require independent verification.
\item The supplied registry labels the statement \texttt{corrected\_verified} and says that it was reconstructed and checked against the official Oberwolfach report. These are supplied verification claims; no source visit or independent verification is asserted here.
\end{itemize}
\end{document}
